% Candidate insertions generated during the May 22 discovery pass.
% This file does not modify last.tex. It collects sections that are ready,
% or close to ready, for later integration.

\subsection{Curvature, knot cost, and the uniqueness of three dimensions}

The same number appears from two independent sides of the geometry.  On the unit sphere \(S^n\), scalar curvature is fixed by dimension:
\[
R_{S^n}=n(n-1).
\]
On the Clifford torus, the adjacent torus knot \(T(n-1,n)\) has Chern-Simons invariant
\[
\mathrm{cs}(T(n-1,n))
=
\frac{((n-1)^2-1)(n^2-1)}{24n(n-1)}
=
\frac{(n-2)(n+1)}{24}.
\]
The inverse invariant \(1/\mathrm{cs}\) is the natural topological cost of maintaining that knot.  Requiring the simplest stable topology in dimension \(n\) to be calibrated against the ambient curvature gives
\[
\frac{1}{\mathrm{cs}(T(n-1,n))}=R_{S^n},
\]
so
\[
\frac{24}{(n-2)(n+1)}=n(n-1),
\]
and therefore
\[
n(n-1)(n-2)(n+1)=24.
\]
For integer \(n>2\), the left side is strictly increasing.  At \(n=3\) it is
\[
3\cdot2\cdot1\cdot4=24,
\]
while the next case already gives
\[
4\cdot3\cdot2\cdot5=120.
\]
Thus \(n=3\) is the unique dimension where intrinsic sphere curvature and the topological cost of the simplest adjacent torus knot use the same natural scale.

This result does not rely on a fitted constant.  It uses the scalar curvature of the unit sphere and the standard Chern-Simons invariant of torus-knot complements.  Its physical reading is direct: in three spatial dimensions, matter as stable topology and gravity as curvature can share a common unit.  In other integer dimensions, the same calibration fails at the arithmetic level.

\subsection{The finite phase sector of the trefoil}

The trefoil complement has fundamental group
\[
\pi_1(S^3\setminus T(2,3))
=
\langle a,b\mid a^2=b^3\rangle.
\]
Taken as a raw knot group, its one-dimensional \(U(1)\) character space is continuous.  Abelianization gives the relation
\[
2a=3b,
\]
and since \(\gcd(2,3)=1\), the abelianization is \(\mathbb Z\).  The six-fold structure appears after imposing the finite closure sector
\[
a^2=b^3=1.
\]
Then \(a\) has two allowed phases and \(b\) has three:
\[
a\in\{\pm1\},
\qquad
b\in\{1,\omega,\omega^2\},
\qquad
\omega=e^{2\pi i/3}.
\]
This gives a finite phase set
\[
\mathbb Z_2\times\mathbb Z_3,
\]
with six characters.  The same \(2\times3\) structure appears in the proton prefactor
\[
\frac{1}{\mathrm{cs}(T(2,3))}=6.
\]
The safe statement is therefore precise: the trefoil group itself has a continuous \(U(1)\) character family, while the closed finite phase sector selected by the relations \(a^2=b^3=1\) has six characters, organized as \(2\times3\).

\subsection{The trefoil as the unique integer torus-knot mass prefactor}

For a torus knot \(T(p,q)\), with \(2\leq p<q\) and \(\gcd(p,q)=1\), the Kirk--Klassen invariant gives
\[
\mathrm{cs}(T(p,q))
=
\frac{(p^2-1)(q^2-1)}{24pq}.
\]
The inverse topological prefactor is therefore
\[
\frac{1}{\mathrm{cs}(T(p,q))}
=
\frac{24pq}{(p^2-1)(q^2-1)}.
\]
For fixed \(p\), this expression decreases as \(q\) increases, so its maximum occurs at the first admissible value \(q=p+1\).  At that value,
\[
\frac{24p(p+1)}{(p^2-1)((p+1)^2-1)}
=
\frac{24}{(p-1)(p+2)}.
\]
For \(p\geq5\), this maximum is at most
\[
\frac{24}{4\cdot7}=\frac67,
\]
which is already below \(1\).  No positive integer inverse prefactor can occur for \(p\geq5\).

Only \(p=2,3,4\) remain.  For \(p=2\),
\[
\frac{1}{\mathrm{cs}(T(2,q))}
=
\frac{16q}{q^2-1}.
\]
Since \(\gcd(q,q^2-1)=1\), integrality forces
\[
q^2-1\mid16.
\]
With \(q>2\), the only admissible solution is \(q=3\), giving
\[
\frac{1}{\mathrm{cs}(T(2,3))}=6.
\]
For \(p=3\), the only cases with value at least \(1\) are \(q=4,5,7,8\), and they give
\[
\frac{12}{5},\quad
\frac{15}{8},\quad
\frac{21}{16},\quad
\frac87,
\]
none of which is an integer.  For \(p=4\), only \(q=5\) remains, giving
\[
\frac43.
\]
Thus \(T(2,3)\) is the unique torus knot whose inverse Chern-Simons invariant is an integer.  The proton prefactor \(6\) is not merely the smallest example; it is the only integer mass prefactor available to torus-knot topology under this invariant.

\subsection{The rational point of the \(S^4\) spectral zeta}

For the conformal scalar on \(S^4\),
\[
\zeta_{\mathrm{conf}}(s)
=
\frac16\sum_{m=1}^{\infty}
\frac{2m+1}{[m(m+1)]^{s-1}}.
\]
At \(s=3\),
\[
\frac{2m+1}{m^2(m+1)^2}
=
\frac{1}{m^2}-\frac{1}{(m+1)^2},
\]
so the sum telescopes:
\[
\zeta_{\mathrm{conf}}(3)=\frac16.
\]
This is the Chern-Simons invariant of the trefoil complement.

The same spectral function also explains why this rational value is special.  Let \(n=s-1\), and write
\[
F_n(m)=\frac{2m+1}{m^n(m+1)^n}.
\]
The partial fractions have coefficients
\[
A_{n,j}
=
(-1)^{n-j}
\frac{j-1}{2n-j-1}
\binom{2n-j-1}{n-j},
\qquad j=2,\ldots,n,
\]
whose absolute values form Catalan-triangle rows.  These coefficients pair the \(1/m^j\) and \(1/(m+1)^j\) terms so that every even zeta value cancels.  For \(n>2\), an odd zeta value always remains: when \(n\) is odd, the highest term \(\zeta(n)\) survives; when \(n\) is even, the highest odd term \(\zeta(n-1)\) survives with coefficient proportional to \(n-2\).

Thus \(s=3\) is the unique integer evaluation where the spectral zeta is purely rational.  The proton sits at the rational closure point of the \(S^4\) spectrum, while higher evaluations generate the odd zeta values that appear in one-loop corrections.

\subsection{The fine-structure running scale}

The bare coupling is
\[
\alpha_{\mathrm{bare}}^{-1}
=
\frac{9\pi^5}{e^3}.
\]
The observed low-energy coupling is obtained by the one-loop running term
\[
\alpha_{\mathrm{phys}}^{-1}
=
\frac{9\pi^5}{e^3}
-
\frac{1}{3\pi}\log\frac94.
\]
The scale ratio \(9/4\) follows from the same projection geometry.  At the bare coupling point, the \(W\)-coordinate is
\[
\cos\sigma_{\mathrm{bare}}=\frac23,
\]
the electromagnetic charge coordinate of the up quark.  If energy scales by inverse projected length, then
\[
\frac{Q}{\mu}
=
\sec\sigma_{\mathrm{bare}}
=
\frac32,
\]
and the one-loop formula uses
\[
\frac{Q^2}{\mu^2}
=
\sec^2\sigma_{\mathrm{bare}}
=
\frac94.
\]
This is the same number written in dimensional form:
\[
\frac94=\left(\frac{3}{2}\right)^2.
\]
The logarithmic correction is therefore the QED running from the \(W=2/3\) geometric scale to the electron scale.

\subsection{The gravitational hierarchy and the \(61\alpha^2\) target}

The Planck/electron mass hierarchy is controlled by the same trefoil and dimension factors:
\[
\ln\frac{m_{\mathrm{Pl}}}{m_e}
\approx
\frac{27\pi^5}{8e^3}
+
\frac{\zeta(5)}{10}.
\]
The first term is
\[
\frac{27\pi^5}{8e^3}
=
\frac38\frac{9\pi^5}{e^3},
\]
where
\[
\frac38
=
\mathrm{cs}(T(2,3))\left(\frac32\right)^2
=
\frac16\cdot\frac94.
\]
The zeta term is three copies of the proton correction,
\[
\frac{\zeta(5)}{10}=3\cdot\frac{\zeta(5)}{30}.
\]
Numerically,
\[
\frac{27\pi^5}{8e^3}+\frac{\zeta(5)}{10}
=
51.52459529358\ldots,
\]
while the value from current constants is
\[
\ln\frac{m_{\mathrm{Pl}}}{m_e}
=
51.52783983692\ldots.
\]
The residual is
\[
0.0032445433385\ldots
=
60.9288414861\ldots\,\alpha^2.
\]
The next correction target is therefore
\[
61\alpha^2.
\]
Adding it leaves a logarithmic residual of about \(3.8\times10^{-6}\), corresponding to a few parts per million in \(G\).  The integer \(61\) is not yet derived; the useful result at this stage is that the remaining correction has been reduced to a specific second-order electromagnetic scale.

\subsection{Heavy-sector multiplicative correction}

The additive zeta correction works across the low and intermediate spectrum, but the \(\pi^7\) electroweak family leaves a coherent residual.  For the heavy sector, the correction is better modeled as a fractional one-loop shift:
\[
\frac{\delta m}{m_{\mathrm{bare}}}
=
\kappa\frac{\alpha}{2\pi},
\qquad
m_{\mathrm{pred}}
=
m_{\mathrm{bare}}\left(1+\kappa\frac{\alpha}{2\pi}\right).
\]
Using the current bare formulas gives
\[
\kappa^\star_t=0.652012858,
\quad
\kappa^\star_W=1.295012119,
\quad
\kappa^\star_Z=1.203927825,
\quad
\kappa^\star_H=1.624995931.
\]
The class structure suggested by these values is
\[
\kappa_t\sim\frac23,
\qquad
\kappa_W\sim\frac43,
\qquad
\kappa_H\sim\frac{13}{8}.
\]
For the \(Z\), the natural form includes weak mixing:
\[
\kappa_Z=\frac43-\lambda\sin^2\theta_W.
\]
With \(\sin^2\theta_W=0.223044624\), the extracted value is
\[
\lambda=0.58017766\ldots.
\]
The small-denominator candidate \(7/12\) is close, but no structural derivation of that denominator is currently known.

With \(\kappa_H=13/8\) and the extracted \(Z\)-mixing coefficient, the full-spectrum geometric mean absolute error is \(0.000013850\%\), and the worst residual is \(0.004443979\%\), carried by the W.  Using small rational values such as \(\lambda=7/12\), \(11/19\), or \(18/31\) for the \(Z\)-mixing coefficient keeps the worst residual the same and leaves the geometric mean in the \(3.6\)--\(4.1\times10^{-5}\%\) range.  The regime switch is strong enough to keep as a discovered correction layer, while the exact origin of the \(\kappa\) values remains a derivation target.
